\section{Grupo fundamental}

\subsection{Arcos y homotopías}
\begin{definicion}[Arco]
  Un arco en el espacio \(X\) es una función continua \(f : I \to X \)
  donde \(I = [0,1]\)
\end{definicion}
%% mariela: q1 ¿todo espacio 1-dimensional es homeomorfo a [0,1]?
%% no dije exactamente eso, un subconjunto convexo y acotado si lo es ,
%% pero debo de preguntar que se refiere
La definición tradicional es ligeramente mas general, cambiando que \(I\)
sea un subconjunto convexo acotado de un espacio unidimensional, pero
siempre podemos recuperar esta definición rescalando la distancia.

\begin{definicion}[Homotopía]
  Dados dos arcos \(f,g : I \to X\), diremos que \(f\) es homotópico a
  \(g\) si existe una función continua \(F : I \times I \to X \) tal que
  \[ \begin{matrix}
      F (x, 0) = f(x), & F (x, 1) = g(x)
     \end{matrix}
  \]
  Donde \(F\) sera llamada una homotopía entre \(f\) y \(g\). La
  existencia de dicha relación entre dos funciones sera denotada por \(f
  \homRelAlt g\).
\end{definicion}
\begin{figure}[h]
  \centering
  \includegraphics[scale=0.3]{./imagenes/homotopia.png}
  \caption{caracterización homotopía}
  \label{fig:homotopia-entre-funciones}
\end{figure}
Podemos pensar en el segundo argumento de una homotopía como el grado de
deformación entre dos funciones.
Si tratamos con arcos \(f,g : I \to X\) que posean los mismos puntos
iniciales y finales, es decir \(f(0) = g(0) = x_0, \; f(1) = g(1) =
x_1 \) podemos definir una relación homotópica ligeramente mas fuerte.

\begin{definicion}[Arco homotopía]
  \(f,g : I \to X\) son \emph{arco homotópicas} entre sí, si tienen los mismos
  puntos inicial y final \(x_0, x_1\) respectivamente y existe una homotopía entre
  ellos tal que cumpla
  \[
    \begin{matrix}
      F(s,0) = f(s), & F(s,1) = g(s) \\
      F(0,t) = x_0,  & F(1,t) = x_1
    \end{matrix}
  \]
  La existencia de esta relación entre dos funciones sera denotada por
  \(f \simeq g\).
\end{definicion}
% todo(slack) hacer algun dibujo de las homotopías entre las curvas
Una vez definida estas relaciones, una pregunta natural es si estas definen
una relación de equivalencia, ya que nos gustaría identificar clases de
curvas como elementos de alguna teoría algebraica. La respuesta es
afirmativa para ambas relaciones (el punto de inicio y final no juegan
un papel relevante). A continuación mostraremos que ``ser homotópicas''
es una relación de equivalencia.
\begin{teorema}
  \(\homRelAlt\) es una relación de equivalencia
\end{teorema}
\begin{proof}
  Hemos de probar que esta relación cumple la reflexividad, simetría y
  transitividad. Sean \(f,g,h : I \to X\) tres arcos arbitrarios.
  \begin{itemize}
  \item La reflexividad es directa pues la función \(F(x,t) = f(x)\) es
    una deformación continua de \(f\) a \(f\), por tanto \(f \homRelAlt
    f\).

  \item La simetría se obtiene de invertir el sentido de deformación de la
    homotopía original. Formalmente, dado \(f \stackrel{.}{\simeq} g\)
    tenemos una homotopía entre estas \((x,t) \mapsto F(x,t)\). A partir
    de aquí podemos definir
    \begin{equation}
      \label{eq:homotopy-simetry}
      (x,t) \mapsto \hat{F}(x,t) := F(x,1-t)
    \end{equation}
    con \(\hat{F}\) una homotopía entre \(g\) y \(f\). Por tanto \(g
    \stackrel{.}{\simeq} f\).

  \item La transitividad se obtiene a partir de de dividir \(I\) en dos
    intervalos donde se deformen individualmente cada homotopía al doble
    del grado. Formalmente dado \(f \stackrel{.}{\simeq} g\) y \(g
    \stackrel{.}{\simeq} h\) representadas por las homotopías \(F\) y
    \(G\) respectivamente, se define
    \[ FG(x,t) = \begin{cases}
        F(x,2t) & t \in [0,\frac{1}{2}] \\
        G(x,2t - 1) & t \in [ \frac{1}{2} , 1]
      \end{cases}
    \]
    Esta es una deformación continua claramente en \((x,t) \in I \times
    [0, \frac{1}{2}) \cup I \times (\frac{1}{2}, 1]\). La continuidad en
    \(I \times \{\frac{1}{2}\}\) proviene de la consistencia en dicho
    punto de ambas homotopías
    \[ F(x,2 \cdot \frac{1}{2}) = g(x) = G(x, 2 \cdot \frac{1}{2} - 1)\]
    lo que nos permite utilizar el \emph{lema del pegamiento}
    \cite[p.~108]{munkres} para obtener la continuidad de \(FG\).
    Obteniendo así \(f \stackrel{.}{\simeq} h\).
  \end{itemize}
\end{proof}

En análisis real, es común trabajar con conjuntos convexos. En arcos
sobre estos conjuntos, veremos una homotopía repetidamente llamada
homotopía de linea recta.
\begin{definicion}[Homotopía de linea recta]\label{def:homotopia-linea}
  Sean \(f,g : I \to X\) dos arcos arbitrarios y sea \(\mathcal C
  \subset X\) un subconjunto convexo. Si \(f(I),g(I) \subset \mathcal
  C\), entonces se define la homotopía de linea recta por
  \[ F(x,t) := (1-t) \cdot f(x) + t \cdot g(x) \]
\end{definicion}
\begin{acotacion}
  En la definición anterior, la continuidad de \(F\) se obtiene de ser
  una combinación convexa de funciones continuas \(f\) y \(g\). De que
  sea una combinación convexa se obtiene la importancia de que existe un
  conjunto \(\mathcal C \subset X\) que contenga a las imágenes.
\end{acotacion}

Con estas definiciones ya podemos empezar a hablar de \([f]\) clases
de equivalencia de funciones bajo una relación homotópica
\[ [f] = \{ g : I \to X \mid f \stackrel{.}{\simeq} g \} \]
Notar que la construcción es valida para las dos relaciones aunque por
nuestra meta de construir el grupo fundamental nos interesa
principalmente la relación arco homotópica, por razones a estudiar mas
adelante.

\subsection{Grupoide de arcos}
Si tenemos dos arcos continuos \(f,g\) tales que el punto final de \(f\)
sea el punto inicial de \(g\) podemos construir un camino continuo que
recorra \(f\) luego recorriendo \(g\). Esta idea es conocida formalmente
como el producto de dos arcos.

\begin{definicion}[Producto de arcos]
Para dos arcos \(f,g : I \to X\) tales que
\(f(1) = g(0)\), se define el producto \(f * g \)
\[ (f*g) (s) \coloneqq \begin{cases}
    f(2s) & s \in [0,\frac{1}{2}] \\
    g(2s - 1) & s \in [\frac{1}{2} , 1]
  \end{cases}
\]
la cual sigue siendo una funcion continua en virtud del lema del
pegamiento.
\end{definicion}

Esta construccion se puede reutilizar para clases
\emph{arco}-homotopicas \([f],[g]\) que compartan punto final e inicial
respectivamente para definir un producto de clases de equivalencia
\[ [f] * [g] \coloneqq [f * g]\]
El cual esta bien definido pues si \(f \simeq f' ,\ g \simeq
g'\) a traves de las homotopias \(F, G\) respectivamente, entonces
\[
  H(s,t) = \begin{cases}
    F(2s,t) & s \in [0, \frac{1}{2}] \\
    G(2s - 1, t) & s \in [\frac{1}{2} , 1]
  \end{cases}
\]
Es la homotopia que relaciona a los arcos que sean homotopicos a
\(f*g\). Esta ademas es continua en virtud otra vez del lema del pegamiento.

\paragraph{} Con esta operacion binaria, una pregunta a hacerse es si
\((\mathcal C (I , X)/\simeq , (*))\) tiene posee estructura de
grupoide, esto equivale a cumplir 3 propiedades
\begin{enumerate}
\item \textbf{Asociatividad} Si \([f] * ([g] * [h])\) esta definido entonces
  tambien lo debe estar \(([f] * [g]) * [h]\) y ademas deben coincidir.
\item \textbf{Identidades izquierda y derecha} Para todo \([f]\) con
  \(x_0, x_1\) puntos inicial y final respectivamente, debe de
existir elementos \([k_{x_0}], [k_{x_1}]\) tales que
\[ \begin{matrix}
    [f] * [k_{x_1}] = [f] & & [k_{x_0}] * [f] = [f]
  \end{matrix}
\]
\item \textbf{Inverso} Para todo \([f]\) clase de arcos con \(x_0, x_1\)
  puntos inicial y final respectivamente debe de existir un elemento
  \([f^{-1}]\) que cumpla
\[ \begin{matrix}
    [f] * [f^{-1}] = [k_{x_0}] & & [f^{-1}] * [f] = [k_{x_1}]
  \end{matrix}
\]
\end{enumerate}
Para probar esto necesitamos primero definir a nuestros candidatos de
\(k_{x_0}, k_{x_1}, f^{-1}\) ademas de algunas funciones auxiliares,
iniciando por los mapeos constantes e identidad en \(I \to I\)
\[ \begin{matrix}
     e_0 : & I \to I & e_0(t) \coloneqq 0 \\
     e_1 : & I \to I & e_1(t) \coloneqq 1 \\
     i :   & I \to I & i(t) \coloneqq t \\
     \bar{i} : & I \to I & i(t) \coloneqq 1 - t
   \end{matrix}
   \]
Ademas, para todo arco \(f : I \to X \), el elemento \(f^{-1} : I \to X \) esta
definido (en el espiritu de \eqref{eq:homotopy-simetry}) por
\[ f^{-1} (s) \coloneqq f (1 - s) \]
Por ultimo, para todo \(x \in X \) se define la curva
constante\footnote{\(k\) por \emph{konst}}
\begin{align*}
  k_x : &I \to X \\
        &t \mapsto x
\end{align*}
Con esto podemos afrontar la demostración
\begin{proof}
Se procedera en orden \(2 \to 3 \to 1\) con lemas en medio. Sea \(f : I
\to X\) el representante de la clase \([f]\). sean \(x_0, x_1\) los
puntos inicial y final de \(f\) respectivamente.

\paragraph{(2).} Dados \(i,e_1\) arcos en \(I\) definidos anteriormente,
es claro que \(i * e_1\) es tambien un arco (continuo) en \(I\); mas aun,
ya que \(I = [0,1]\) es convexo, se tiene \(i * e_1 \simeq i\) con la
homotopia de la linea recta entre estos denotada por \(C\).

Tambien conocemos que la composición de funciones continuas es
continua, de lo cual podemos decir que dada una homotopia \(H : I \times
I \to I\) la composicion \( f \circ H : I \times I \to X\) es una homotopia.

Por otro lado, necesitamos conocer como se comporta la composicion con
respecto a nuestro producto \((*)\), para esto tenemos el siguiente lema
\begin{lema}[Distributividad de la composicion sobre producto]
\label{lema:dist-composicion-producto}
\[\forall a,b : I \to I,\ \forall f : I \to X,\ f \circ (a * b) = (f
\circ a) * (f \circ b) \]
\end{lema}
\begin{proof}
  \[ f \circ (a*b) (s) =
    \begin{cases}
      f (a(2s)) & s \in [0,\frac{1}{2}] \\
      f (b(2s - 1)) & s \in [\frac{1}{2} , 1]
    \end{cases}
    = (f \circ a) * (f \circ b) (s)
  \]
\end{proof}

Luego esto nos permite afirmar en especifico que
\[ f \circ i \simeq f \circ (i * e_1) = (f \circ i) * (f \circ e_1) \]
en virtud de la homotopia
\[ f \circ C : I \times I \to X \]
Por tanto
\begin{equation}\label{eq:homequiv2.1}
[f \circ i] = [f \circ (i * e_1)] = [(f \circ i) * (f \circ e_1)]
\end{equation}
Si tomamos este ultimo termino, sabemos que por definicion
\begin{equation}\label{eq:homequiv2.2}
[(f \circ i) * (f \circ e_1)] = [(f \circ i)] * [(f \circ e_1)]
\end{equation}
y es claro tambien que
\[
  \begin{matrix}
    f \circ i = f & f \circ e_1 = k_{x_1} \\
  \end{matrix}
\]
\begin{equation}\label{eq:homequiv2.3}
[(f \circ i)] * [(f \circ e_1)] = [f] * [k_{x_1}]
\end{equation}
juntando
\eqref{eq:homequiv2.1},\eqref{eq:homequiv2.2},\eqref{eq:homequiv2.3} nos
muestra que \([k_{x_1}]\) es la identidad por la derecha. Podemos
hacer un proceso analogo para \([k_{x_0}]\) como identidad por la
izquierda, asi probando (2).

\paragraph{(3).} De manera similar a la pregunta anterior, notar que \(i,
\bar{i}\) definidos anteriormente cumplen
\[ i * \bar{i} \simeq e_0 \]
por la homotopia de linea recta en el espacio \(I\) que es convexo.
Utilizando el lemma \eqref{lema:dist-composicion-producto} tenemos que
\[ [f * f^{-1}] = [f \circ (i * \bar{i})] = [f \circ e_0] = [k_{x_0}] \]
\[ \implies [f] * [f^{-1}] = [k_{x_0}] \]
De manera analoga se prueba la inversa por la izquierda.

\paragraph{(1).}
Por definicion, tenemos que las diferentes asociaciones nos resultan en
diferentes representates de clase
\[
  \begin{matrix}
    (f * g) * h \, (t) =
    \begin{cases}
      f (4t), & t \in [0, \frac 1 4] \\
      g (4t - 1), & t \in [\frac 1 4, \frac 1 2] \\
      h (2t - 1), & t \in [\frac 1 2, 1] \\
    \end{cases} &
    f * (g * h) \, (t) =
    \begin{cases}
      f (2t), & t \in [0, \frac 1 2] \\
      g (4t - 2), & t \in [\frac 1 2, \frac 3 4] \\
      h (4t - 3), & t \in [\frac 3 4, 1] \\
    \end{cases}
  \end{matrix}
\]
Sea \(T : I \to I\) definida por
\[ T (t) :=
  \begin{cases}
    2t & t \in [1 , \frac 1 4] \\
    t + \frac 1 4 & t \in [\frac 1 4, \frac 1 2] \\
    \frac t 2 + \frac 1 2 & t \in [\frac 1 2, 1]
  \end{cases}
\]
\(T\) es una funcion continua por el lema del pegamiento. Por ser \(I\)
un dominio convexo, se tiene que \(i \simeq T\), luego
\[ [(f * g) * h] = [\big( (f * g) * h \big) \circ i] = [\big( (f * g) *
  h \big) \circ T ] = [f * (g * h)] \]
\end{proof}


\subsection{Grupo fundamental}
Sobre nuestros arcos tenemos una estructura de grupoide porque no siempre
podemos hacer coincidir los puntos inicial y final. Si nos fijaramos en
un solo punto \(x_0 \in X\) tal que todas los arcos \(f : I \to X\)
cumplieran que
\[ f(0) = x_0 = f (1) \]
Entonces las propiedades que demostramos anteriormente nos dirian que el
espacio \((X / \simeq_{ah}, *)\) es un grupo.
\begin{definicion}[Grupo fundamental]
  Sea \(X\) un espacio topologico y \(x_0 \in X\) un punto. La coleccion
  \[ \pi_1 (X,x_0) = \{ [f] \mid f : I \to X, f \text{ continua}\}\]
  \[ [f] := \{ g : I \to X \mid f \simeq_{ah} g\}\]
  junto con el producto de caminos \((*)\) forman el llamado
  \textbf{grupo fundamental} \((\pi_1(X,x_0), (*))\)
\end{definicion}
Nuestro interes es mostrar que para un espacio dado, su grupo fundamental
es una invariante topologica. Esto quiere decir que dado dos espacios puntuados
\((X,x_0)\) e \((Y,y_0)\) homeomorfos, es decir que existe un isomorfismos continuo
entre ellos, debe de existir un respectivo isomorfismo de grupos entre
\(\pi_1(X,x_0)\) y \(\pi_1(Y,y_0)\). Para esto debemos recordar que
significa ser isomorfo entre grupos
\begin{definicion}
  Sean \((G,(*)), (G', (\cdot))\) does grupos. \(F : G \to G'\) es un
  homomorfismo entre grupos si y solo si
  \[ \forall x,y \in G,\ F (x * y) = F(x) \cdot F(y)\]
  \[ e \text{ identidad en } G \implies F(e) \text { identidad en } G' \]
\end{definicion}

\paragraph{Ejemplo \(\mathbb{R}^2 - \{(1,1)\}\).} Este espacio es
relativamente sencillo de visualizar, haremos un argumento intuitivo para
ver su grupo fundamental. Sea \(f : I \to \mathbb{R}^2 - \{(1,1)\}\) un
camino arbitrario

%% old
Una pregunta natural ha
hacerse es si la eleccion de punto influye en el tipo de grupo que se
obtiene en \((X / \simeq_{ah}, *)\). En general la respuesta es
afirmativa, lo cual complica nuestra tarea puesto que queremos definir
una invariante topologica que dependa del espacio total. Pero podemos
recuperar nuestro objetivo si nos restringimos a espacios \(X\) que sean
arco-conexos.
\begin{definicion}[Arco-conexidad]
  Diremos que el espacio \(X\) es \emph{arco conexo} si para todo par de
  puntos \(x,y \in X\) existe un arco continuo \(f : I \to X\) tal que
  cumpla
  \[ f(0) = x \qquad f(1) = y \]
\end{definicion}
Para estos espacios, denotamos a \((X / \simeq_{ah}, *)\) como
\((\pi_1(X,x_0), *)\) y le llamaremos \textbf{grupo fundamental} de \(X\) con
respecto a \(x_0\). Es claro que la importancia del punto \(x_0\) es poca
considerando la arco-conexidad. Puesto que si tengo \([f] \in
\pi_1(X,x_0)\) y \(\alpha : I \to X\) un arco entre \(x_0\) y \(x_1\),
existe un camino \([\hat{\alpha} * f * \alpha] \in \pi_1(X,x_1)\), es
decir un isomorfismo entre \(\pi_1 (X, x_0)\) y \(\pi_1 (X, x_1)\). Aun
asi es conveniente reificar los puntos bases de los grupo puesto que no
necesariamente existe un camino canonico entre ellos.
% La razon de porque aun guardamos a \(x_0\) como
% informacion del grupo es que a priori los isomorfismos entre \(\pi_1 (X,
% x_0)\) y \(\pi_1 (X, x_1)\) pueden depender del arco \(\alpha\) elegido,
% de los cuales ningun

%todo(slack): ejemplo de R^2 sin el punto
%todo(slack): ver la discucion de la propiedades functoriales (talvez?)